Similarly, every $\type{T}\in\KFpt(\ell)$ is contained
		in $F(\type{T})$, and the construction agrees with $f$ on $S$.
		We check the weak successor condition of
		Definition~\ref{def:shape-pres}.
		Consider a successor
		$\type{Q}=\S(\type{T},\bar p,\type{C})$
		with $\ell(\type{T})=n$.
		If $n<\ell-1$, all involved levels are fixed and the successor
		is unchanged.
		If $n=\ell-1$, the base and its parameters are fixed, while the
		old successor is an initial segment of $F(\type{Q})$;
		the inserted coordinate creates precisely the permitted extra
		level. In particular, equality with $F(\type{Q})$ must
		\emph{not} be asserted at this boundary.
		If $n\geq\ell$, the old ordinary column and the terminal
		$L^+$-letter of the successor are copied by insertion.
		Its empty or singleton canonical parameter is transported
		by the same construction, using the first missing $E$-coordinate.
		The uniqueness of the successor decomposition then gives
		$\S(F(\type{T}),F(\bar p),\type{C})=F(\type{Q})$.
		These cases give the required weak successor property; roots
		are fixed since $\ell>0$, and the level map is the strictly
		increasing function fixing levels below $\ell$ and adding one
		from $\ell$ onward.
		Consequently $F\in\M$ and it skips precisely the level $\ell$.
